\subsection*{Croqtile 函数}

先看函数签名。普通 C++ 函数的声明主要约定类型，Croqtile 还把数据形状写在参数和返回值中。这样，读者可以在进入函数体之前，先了解输入内容与预期的结果形状。

\begin{lstlisting}[language=croqtile,caption={核函数签名}]
__co__ s32 [4, 8] ele_add(s32 [4, 8] lhs, s32 [4, 8] rhs) {
\end{lstlisting}

前缀 \co@__co__@ 标记 Croqtile 函数。\co@s32 [4, 8]@ 把 32 位有符号整数元素与二维形状组合起来，表示四行八列。输入与返回值都采用 \co@type [shape] name@ 的形式。编译器据此检查秩并推导索引范围；依赖调用方实际尺寸的约束，还可能需要运行时检查。

常见元素类型包括 \co@s32@、\co@f16@、\co@bf16@、\co@f32@，以及 \co@f8_e4m3@ 和 \co@f8_e5m2@。本例使用 \co@s32@，便于观察逐元素加法的结果。

\begin{lstlisting}[language=croqtile,caption={分配输出}]
s32 [lhs.span] output;
\end{lstlisting}

这一行声明输出张量。\co@lhs.span@ 取出输入的完整形状，因此 \co@output@ 也具有 \co@[4, 8]@ 的形状。修改输入形状时，这条声明的输出形状会随之改变；本例中固定的函数签名和 \co@parallel@ 范围仍需同步调整。

\begin{lstlisting}[language=croqtile,caption={并行函数体}]
parallel {i, j} by [4, 8]
  output.at(i, j) = lhs.at(i, j) + rhs.at(i, j);
\end{lstlisting}

语句 \co@parallel {i, j} by [4, 8]@ 声明 32 个独立实例，每个实例负责一个元素位置。例如 \co@(i, j) = (1, 3)@ 的实例读取两份输入的这一位置，再把和写入 \co@output@ 的相同位置。不同实例的写入位置不重叠，也不依赖其他实例产生的结果，因此这些工作可以并行执行。

没有空间说明符时，由编译器选择实例到执行资源的映射。语句确定工作彼此独立，并不要求每个元素都对应一个物理线程，也不要求所有实例同时运行。

\begin{lstlisting}[language=croqtile,caption={返回结果}]
return output;
\end{lstlisting}

返回结果的形状与元素类型必须符合函数签名。主机侧收到的 \co@croq::spanned_data@ 拥有结果缓冲区，并携带形状信息，可通过 \co@[i][j]@ 访问元素。
